%% =====================================================================
%%  Page 1 -- exactly four items. Nothing else.
%%  Team number / problem letter / title / summary (<=300 words).
%%  No names, no school, no date, no logo, no abstract heading other than
%%  "Summary".
%%
%%  NOTE: the scenario and the numbers below are a worked example that is
%%  consistent with code/blowdown.py. Replace them with your own problem.
%% =====================================================================
\begin{titlepage}
%%  Layout after the MCM/ICM summary sheet (github.com/latexstudio-org/mcmthesis):
%%  the two identifiers sit in opposite corners of a ruled band, the title takes
%%  the optical centre, the summary follows. No running head on this page, so
%%  nothing competes with the four items.
\thispagestyle{empty}
\vspace*{0.6cm}
\noindent
\begin{minipage}[t]{0.45\textwidth}
  \raggedright
  {\scshape Problem}\\[0.35em]{\LARGE\bfseries\problemltr}
\end{minipage}\hfill
\begin{minipage}[t]{0.45\textwidth}
  \raggedleft
  {\scshape Team}\\[0.35em]{\LARGE\bfseries\teamnum}
\end{minipage}\par

\vspace{1.1em}
\noindent\rule{\textwidth}{0.6pt}\par
\vspace{1.7em}

{\centering
  {\LARGE\bfseries\setstretch{1.12}\papertitle\par}
  \vspace{1.5em}
  % 64 of the 74 gold papers call this block "Abstract"; the rule calls it
  % "a summary ... which will be an abstract", so either word is fine.
  % Change it back to Summary here if the team prefers the rule's wording.
  {\large\scshape Abstract}\par
}
\vspace{0.9em}

%%  A 200-300 word block at full width reads as a wall; the small inset gives
%%  it a measure of about 80 characters and marks it as a separate element.
{\setlength{\leftskip}{0.05\textwidth}%
 \setlength{\rightskip}{0.05\textwidth}%
 \setlength{\parindent}{0pt}%
 \setstretch{1.04}%
A penalty kick is decided before the ball has travelled three metres, and the
goalkeeper's side is committed about $0.1$ s before the kick. We ask whether the
goal-plane crossing point can be recovered from two measured positions, at
$0.05$ s and $0.10$ s after the kick. A six-parameter transverse-force model,
comprising a launch velocity, a drag coefficient and two transverse
coefficients, is fitted to the six coordinates by our own damped Gauss-Newton
routine at an assumed noise of $5$ mm. The fitted trajectory crosses the plane
at $y=2.785$ m, $z=1.945$ m, $t=0.399$ s, $0.87$ m inside the post and $0.50$ m
below the crossbar, with $90$ per cent intervals $y\in[2.455,3.121]$ m and
$z\in[1.617,2.276]$ m; the whole ball passes inside the frame with probability
$0.971$. The four complete models place the point between $2.549$ and
$2.785$ m, and the model envelope of calibrated intervals is
$y\in[2.125,3.121]$ m.

Imposing extreme assumptions widens this to a scenario envelope,
$y\in[2.259,3.935]$ m and $z\in[1.169,2.698]$ m, which is a stress range and not
a confidence interval. Coverage is tested on the estimator itself, with truths
from the same law and from a law containing an extra vertical force, and the
answer is unchanged when common-mode offsets or time-stamp errors are admitted.
The method
is also a design statement: the two epochs of the problem give
$\sigma_y=211$ mm, while fixes at $0.11$ s and $0.31$ s give $12$ mm, a factor of
$18$. The prediction is ready $0.11$ s after the kick, $135$ ms before the
goalkeeper can resolve height and $0.29$ s before the ball arrives.\par}

\ifkeywords
  \vspace{1.0em}
  \begingroup
  \setlength{\leftskip}{0.05\textwidth}%
  \setlength{\rightskip}{0.05\textwidth}%
  \setlength{\parindent}{0pt}%
  \noindent\textbf{Keywords:}~\keywords\par
  \endgroup
\fi

\vfill
\end{titlepage}
\setcounter{page}{2}% titlepage resets the page counter; keep numbering continuous
